\documentclass[10pt,twocolumn]{article}
\usepackage[margin=0.75in]{geometry}
\usepackage{amsmath,amssymb,amsfonts}
\usepackage{graphicx}
\usepackage{booktabs}
\usepackage{xcolor}
\usepackage{titlesec}
\usepackage[hidelinks]{hyperref}
\usepackage{caption}
\captionsetup{font=small,labelfont=bf}
\titleformat{\section}{\normalfont\large\bfseries}{\thesection}{0.6em}{}
\titleformat{\subsection}{\normalfont\bfseries}{\thesubsection}{0.5em}{}
\setlength{\parskip}{2pt}

\title{\vspace{-2em}\textbf{One State to Draw Them All}\\[0.2em]
{\large A Unified Latent State for Expression, Space, Morphology, and Time in Single-Cell Biology}}
\author{Yuchan Lee\\{\small Built with Claude --- Life Sciences Hackathon (Researcher Track)}}
\date{}

\begin{document}
\maketitle

\begin{abstract}
\small
Single-cell measurements are fragmented across modalities: transcriptomics,
morphology, spatial context, and temporal dynamics are acquired on different
platforms with no shared coordinate system. We ask whether a \emph{single}
128-dimensional latent state $S$ can serve as that coordinate system. We train
one shared encoder $E:\mathbb{R}^{422}\!\to\!\mathbb{R}^{128}$ on 200{,}000
cells from a 10x Xenium human colon-cancer sample, with three heads that decode
$S$ back to (i) expression, (ii) a self-supervised morphology embedding, and
(iii) spatial neighborhood structure. A classifier-free-guidance (CFG)
diffusion decoder then \emph{generates} the cell's microscopy image directly
from $S$, and the \emph{frozen} encoder transfers an independent EMT time-course
into the same space. Every axis is validated against a shuffle/identity control:
expression $R^2{=}0.574$, spatial $R^2{=}0.331$, morphology gain $+0.30$, time
Spearman $\rho{=}0.50$, and a generation cell-specificity gap of $+0.17$ at
$w{=}3$. We report the full method, the equations, and an honest account of the
failure (blurry hallucination) that motivated the CFG fix. Code, weights, and a
public interactive demo are released.
\end{abstract}

\section{Introduction \& Motivation}
A ``world model'' for cells should predict how a cell looks, where it sits, and
how it changes---all from a compact internal state. In practice, each of those
axes lives in a separate dataset with an incompatible representation. A
transcriptomic latent (e.g.\ scVI) knows nothing about pixels; a morphology
embedding (e.g.\ DINOv2 on a DAPI crop) knows nothing about neighbors; a
time-course lives on yet another platform. The central hypothesis of this work
is that these are \emph{views of one underlying state}, and that a single
$S$ learned to be simultaneously decodable to all of them is both (a) a useful
compression and (b) a bridge that lets an axis missing in one dataset be
\emph{reconstructed} from another.

Our contribution is not a new primitive. CFG, DINOv2, GNNs, and FiLM
conditioning are all established. The contribution is the \emph{composition}:
four axes fused into one predictable state, trained end-to-end, with a negative
control at \emph{every} step---and a running, servable system rather than a
schematic.

\begin{figure*}[t]
\centering
\includegraphics[width=0.92\textwidth]{unified_architecture_transformer.png}
\caption{\textbf{Architecture.} A shared encoder maps 422-gene expression to a
128-d state $S$ (left). Three auxiliary heads decode $S$ to expression,
morphology (DINOv2-384), and spatial neighborhood; a CFG diffusion U-Net (right)
generates the cell image from $S$ via FiLM conditioning in every ResBlock. Time
is a \emph{transfer}: the frozen encoder projects an independent EMT time-course
into the same $S$.}
\label{fig:arch}
\end{figure*}

\begin{figure*}[t]
\centering
\includegraphics[width=0.98\textwidth]{unified_worldmodel_demo.png}
\caption{\textbf{The system at a glance.} One shared 128-d state $S$ predicts all
four axes (top). Real Xenium cells vs.\ images generated purely from $S$
(middle-left); the EMT time trajectory in $S$ (middle-right); and the four-axis
quantitative summary, every axis beating its shuffle control (bottom). A live
public demo lets anyone reproduce the middle-left panel interactively.}
\label{fig:hero}
\end{figure*}

\section{Related Work}
\textbf{Expression$\to$image generation.} GeneFlow (2025) couples an
attention-based RNA encoder with a conditional U-Net trained by rectified flow
and reports FID 20.73 on the same Xenium platform. Spatia (2025) scales to
49 donors, 17 tissue types and 12 disease states ($\sim$17M cell--gene training
pairs), fusing morphology, expression, and space. Both exceed us on
expression$\to$image \emph{alone}, in scale and absolute image metrics.
\textbf{Single-cell latents.} scVI provides a probabilistic negative-binomial
latent; our encoder is deterministic and, crucially, shared across modalities.
\textbf{Our niche.} The four-axis fusion into one state with a control at every
step, including a cross-platform \emph{time transfer}, is what we did not find in
prior work.

\subsection{Where we sit in the Virtual Cell landscape}
The ``virtual cell'' has become one of the most active frontiers in
computational biology. The Arc Institute's STATE model was trained on
observational data from 167M cells and perturbational data from over 100M cells
across 70 human cell contexts; its Virtual Cell Challenge (\emph{Cell}, 2025)
frames the central task as predicting a cell's transcriptomic response to a
genetic perturbation in a held-out cell type. CZI's TranscriptFormer spans 112M
cells across 12 species, and its rBio distills such a model into a
conversational reasoner. The common thread is decisive: these are, almost
without exception, \textbf{single-modality (transcriptomic) models whose goal is
perturbation$\to$expression prediction}. They embed expression in a latent space
and attend over genes; they do not draw the cell, place it in tissue, or move it
through time.

Our work is deliberately orthogonal (Table~\ref{tab:vc}). We do not compete on
transcriptomic scale---we have $2\times10^5$ cells against their $10^8$. We ask a
different question: can one latent state be \emph{simultaneously} an expression
code, a morphology code, a spatial code, and a temporal coordinate? Every
flagship virtual-cell model answers only the first. Two recent spatial-image
works---GeneFlow and Spatia---do generate morphology from expression and exceed
us there, but neither unifies all four axes into a single predictable state with
a negative control at each one, nor demonstrates a cross-platform temporal
transfer. That specific combination is our contribution.

\begin{table*}[t]\centering\small
\caption{Positioning against representative virtual-cell and spatial-generative
models. ``Axes'' = which of \{expression E, morphology M, space P, time T\} the
single model represents. Scale figures are as reported by each work.}
\label{tab:vc}
\begin{tabular}{@{}lccccc@{}}
\toprule
Model & Primary task & Modalities / axes & Generates image? & Scale & Negative control?\\
\midrule
Arc STATE & perturbation$\to$expression & E & no & 167M obs / 100M pert & benchmark\\
CZI TranscriptFormer & cross-species expression & E & no & 112M cells, 12 sp. & benchmark\\
scGPT & multi-omics / perturbation & E & no & 33M cells & task metrics\\
GeneFlow (2025) & expression$\to$image & E, M & \textbf{yes} (FID 20.73) & Xenium (1 platform) & FID\\
Spatia (2025) & expression+space+morph & E, M, P & \textbf{yes} & $\sim$17M cell--gene pairs & scale\\
\textbf{Ours} & \textbf{unified state, all axes} & \textbf{E, M, P, T} & \textbf{yes} (gap +0.17) & 200k cells & \textbf{every axis}\\
\bottomrule
\end{tabular}
\end{table*}

\section{Method}
\subsection{Preprocessing}
Raw transcript counts $c\in\mathbb{Z}_{\ge0}^{422}$ per cell are normalized to a
common scale, log-transformed, and standardized per gene:
\begin{equation}
x_g \;=\; \frac{\log\!\big(1 + m\,c_g / \textstyle\sum_{g'} c_{g'}\big) - \mu_g}{\sigma_g},
\end{equation}
where $m$ is the dataset median total count ($m{=}64$), and $(\mu_g,\sigma_g)$
are the per-gene mean and standard deviation. The same $(\mu_g,\sigma_g,m)$ are
stored and reused verbatim for the time-transfer experiment.

\subsection{Shared encoder and auxiliary heads}
The encoder is a multilayer perceptron with two GELU/LayerNorm blocks,
\begin{equation}
S \;=\; E(x) \;=\; W_3\,\phi\big(W_2\,\phi(W_1 x)\big)\in\mathbb{R}^{128},
\end{equation}
$\phi(\cdot)=\mathrm{LN}(\mathrm{GELU}(\cdot))$, with widths $422\!\to\!512\!\to\!256\!\to\!128$.
Three heads $h_k(S)=U_k\,\mathrm{GELU}(V_k S)$ decode $S$:

\noindent\textbf{(1) Expression reconstruction.} $\hat{x}=h_{\text{expr}}(S)$,
loss $\mathcal{L}_{\text{recon}}=\lVert \hat{x}-x\rVert_2^2$.

\noindent\textbf{(2) Morphology.} The target is a self-supervised embedding of
the cell's $64{\times}64$ DAPI crop $I$ obtained from a frozen DINOv2 ViT-S/14,
$z=\mathrm{DINOv2}(I)\in\mathbb{R}^{384}$. The head regresses it:
$\mathcal{L}_{\text{morph}}=\lVert h_{\text{morph}}(S)-z\rVert_2^2$. Using a
self-supervised target (rather than raw pixels) isolates \emph{shape/texture}
statistics from background intensity.

\noindent\textbf{(3) Spatial GNN.} Build a $k$-nearest-neighbor graph
($k{=}12$) on cell centroids $(x,y)$. For center cell $i$ with neighbors
$\mathcal{N}(i)$, a message-passing layer predicts the center state from the
mean neighbor state:
\begin{equation}
\hat{S}_i \;=\; g\Big(\tfrac{1}{|\mathcal N(i)|}\!\!\sum_{j\in\mathcal N(i)}\! S_j\Big),\quad
\mathcal{L}_{\text{spatial}}=\lVert \hat S_i - S_i\rVert_2^2 .
\end{equation}
The encoder and all heads are trained jointly:
\begin{equation}
\mathcal{L}_{\text{joint}}=\mathcal{L}_{\text{recon}}+\mathcal{L}_{\text{morph}}+\tfrac12\mathcal{L}_{\text{spatial}}.
\end{equation}

\subsection{S-conditioned CFG diffusion decoder}
To \emph{generate} the microscopy image from state, we train a denoising
diffusion model conditioned on $S$. With a cosine noise schedule
$\{\beta_t\}_{t=1}^{T}$ ($T{=}400$) and
$\bar\alpha_t=\prod_{s\le t}(1-\beta_s)$, the forward process is
\begin{equation}
q(x_t\mid x_0)=\mathcal N\!\big(\sqrt{\bar\alpha_t}\,x_0,\;(1-\bar\alpha_t)\mathbf I\big).
\end{equation}
A U-Net $\epsilon_\theta(x_t,t,S)$ predicts the noise. Conditioning enters
through a vector $c=\mathrm{MLP}_t(\text{emb}(t))+\mathrm{MLP}_S(S)$, injected by
\textbf{FiLM} in \emph{every} ResBlock at every resolution:
\begin{equation}
\mathrm{FiLM}(h,c)=h\odot(1+\gamma)+\beta,\qquad [\gamma,\beta]=W_{\!f}\,c .
\end{equation}
For classifier-free guidance we drop $S$ to a learned null token $\varnothing$
with probability $p{=}0.1$ during training, minimizing
$\mathbb E_{t,x_0,\epsilon}\lVert\epsilon-\epsilon_\theta(x_t,t,S_{\text{drop}})\rVert_2^2$.
At sampling time we extrapolate:
\begin{equation}
\tilde\epsilon = \epsilon_\theta(x_t,t,\varnothing) + w\big(\epsilon_\theta(x_t,t,S)-\epsilon_\theta(x_t,t,\varnothing)\big),
\end{equation}
and denoise with 50 DDIM steps,
$x_{t-1}=\sqrt{\bar\alpha_{t-1}}\,\hat x_0+\sqrt{1-\bar\alpha_{t-1}}\,\tilde\epsilon$,
$\hat x_0=(x_t-\sqrt{1-\bar\alpha_t}\,\tilde\epsilon)/\sqrt{\bar\alpha_t}$.
Weights are tracked with an exponential moving average (decay $0.999$).

\subsection{The cell-specificity gap (why not FID)}
A decoder can look sharp yet ignore its condition---generating a plausible but
\emph{generic} cell (``hallucination''). FID cannot detect this. We therefore
measure whether a generated image matches \emph{its own} target more than a
random one. With generations $\hat I_i$ and reals $I_i$, and a permutation
$\pi$,
\begin{equation}
\mathrm{gap} \;=\; \underbrace{\tfrac1n\textstyle\sum_i\lVert \hat I_i-I_{\pi(i)}\rVert^2}_{\text{MSE to shuffled}} \;-\; \underbrace{\tfrac1n\textstyle\sum_i\lVert \hat I_i-I_i\rVert^2}_{\text{MSE to true}} .
\end{equation}
$\mathrm{gap}>0$ means generation is cell-specific; $\mathrm{gap}\approx0$ is
exactly the hallucination signature we diagnosed and fixed (\S\ref{sec:history}).

\subsection{Time as a transfer, not a fourth training axis}
Xenium provides expression, morphology, and space at single-cell resolution but
is a static snapshot. To test whether $S$ carries temporal structure, we take an
\emph{independent} EMT scRNA-seq time-course (A549 under TGF-$\beta$1,
$t\in\{0,8\text{h},1\text{d},3\text{d},7\text{d}\}$), keep the 175 genes shared
with the Xenium panel, zero-pad to 422, apply the \emph{stored} Xenium
normalization, and project with the \emph{frozen} encoder: $S^{\text{EMT}}=E(x^{\text{EMT}})$.
No retraining. We then measure the correlation between real time and the
projection of $S^{\text{EMT}}$ onto the principal temporal axis.

\section{Development History: What We Actually Tried}
\label{sec:history}
This project was not a straight line. The path itself is evidence of how deeply
the design space was probed, and each turn was forced by a concrete negative
result rather than a guess.

\noindent\textbf{(1) A na\"ive time-series world model felt biologically thin.}
Our first system predicted how a perturbed cell-state \emph{distribution} shifts
over time with an optimal-transport transition model (Sinkhorn $W_2$), beating an
identity baseline by $+21.5\%$. It worked---but we judged the biological
significance to be weak: predicting a distribution drift is not the same as
understanding a cell. This pushed us toward \emph{pixels}.

\noindent\textbf{(2) A deterministic decoder produced correct-but-blurry blobs.}
Regressing the microscopy image from expression with an MSE decoder beat the mean
image by $+22\%$ and a shuffled control by $+41\%$---so it captured
\emph{something} real---but every output was a soft ellipse. MSE on a
many-to-one map (one expression profile, many valid shapes) provably collapses to
the conditional mean. We measured this directly: reconstructed variance was only
$23\%$ of real. The blur was not a bug to tune away; it was the mathematically
correct answer to the wrong objective.

\noindent\textbf{(3) A diffusion decoder was sharp but \emph{hallucinated}.}
Switching to diffusion made images sharp (sharpness $0.05\!>\!$ real $0.032$),
but we caught a subtle failure with our own control: MSE-to-true ($0.100$)
\emph{equaled} MSE-to-shuffled ($0.099$). The decoder was drawing beautiful,
generic cells that ignored the condition. This is precisely why we do not trust
FID alone, and why we defined the cell-specificity gap.

\begin{figure}[t]\centering
\includegraphics[width=\linewidth]{sharp_vs_blob_comparison.png}
\caption{The failure that shaped the design. Real cells (top); the MSE decoder
is cell-specific but blurry (middle); plain diffusion is sharp but hallucinated
(bottom)---MSE-to-true equals MSE-to-shuffled. This complementary failure
motivated CFG.}\label{fig:blob}
\end{figure}

\noindent\textbf{(4) CFG closed the loop.} Classifier-free guidance with FiLM
conditioning at every resolution restored cell-specificity: the gap grew
monotonically with guidance weight, reaching $+0.17$ at $w{=}3$ while staying
sharp. The condition finally mattered.

\noindent\textbf{(5) From five separate models to one shared state.} Each axis
above had lived in its own network. The final step---and the point of the
paper---was to force them through a \emph{single} $S$: one encoder, jointly
decodable to expression, morphology, and space, then generating images from that
same $S$, and finally absorbing time by transfer. Along the way we also solved a
mundane but decisive engineering problem: uploading data to the GPU was
$100\times$ slower than having the GPU pull the 13\,GB source directly from the
public CDN ($20$\,MB/s), which is what made 200k-cell training feasible at all.

\section{Experimental Setup}
\subsection{Compute environment}
All training was performed on a single \textbf{NVIDIA A100-80GB} GPU rented
on-demand through Modal, driven programmatically from Claude Science. The
container image was CUDA 12.4.1 with Python 3.11, PyTorch 2.x, and
mixed-precision (AMP) enabled. Data persisted on a Modal Volume
(\texttt{xenium-unified}) so that every job---training, checkpointing,
sampling---read and wrote the same 1.3\,GB preprocessed tensor without
re-uploading. Local orchestration ran CPU-only (14 cores, 64\,GB RAM); no
figure or number in this report required a local GPU.

\subsection{Data pipeline}
The source is the 10x Genomics Xenium
\texttt{Human\_Colon\_Cancer\_P1\_CRC\_Add\_on\_FFPE} bundle (307{,}762 cells,
422-gene panel, a 6.5\,GB morphology OME-TIFF). Naively uploading the extracted
tensors from the local machine to the GPU ran at $\sim0.2$\,MB/s. We instead had
the GPU container \emph{pull the 13\,GB source zip directly from the public 10x
CDN} at $20.8$\,MB/s---a $100\times$ speedup that made 200k-cell preprocessing a
4-minute job rather than a multi-hour bottleneck. For each of $N{=}200{,}000$
cells we range-extracted three zip members and produced: (i) a 422-dim
expression vector, (ii) a $64{\times}64$ DAPI morphology patch centered on the
cell, (iii) its $(x,y)$ centroid, and (iv) a DINOv2 ViT-S/14 embedding of the
patch. The JPEG2000-compressed OME-TIFF required the \texttt{imagecodecs}
decoder---a concrete pitfall we hit and record for reproducibility.

\subsection{Training configuration}
\textbf{Unified encoder + heads.} AdamW ($\text{lr}{=}3\!\times\!10^{-4}$,
weight decay $10^{-5}$), cosine-annealed over 100 epochs, batch size 1024, loss
weights $(\lambda_{\text{recon}},\lambda_{\text{morph}},\lambda_{\text{spatial}})
=(1,1,0.5)$, spatial graph $k{=}12$. Total 0.72M parameters; wall-clock
$\approx2$ minutes.

\textbf{S-conditioned diffusion.} AdamW ($3\!\times\!10^{-4}$, weight decay
$10^{-4}$), cosine schedule, batch size 256, $T{=}400$ cosine-$\beta$ steps,
CFG drop $p{=}0.1$, EMA decay $0.999$, U-Net base width 96, conditioning
dimension 256 ($\sim$7M parameters). Sampling uses 50 DDIM steps. To survive
pre-emption, the job committed an EMA checkpoint and a progress JSON to the
Volume every few epochs; the run converged (loss $\approx0.0233$) and was
stopped at epoch 80/120 at $\sim$103\,s/epoch, with all reported figures drawn
from the epoch-80 EMA weights.

\textbf{Time transfer.} No training: the frozen encoder projected 3{,}133
main-arm EMT cells (GSE147405, A549 TGF-$\beta$1) using the 175 shared genes and
the stored Xenium normalization.

\begin{figure}[t]\centering
\includegraphics[width=\linewidth]{unified_training.png}
\caption{Joint 3-axis training of the shared state (200k cells). Left:
simultaneous per-task loss. Right: on held-out cells, $S$ beats its shuffle
control on all three axes.}\label{fig:train}
\end{figure}

\subsection{Reproducibility}
The four numbered scripts (\texttt{01\_data\_prep} $\to$
\texttt{04\_sample}) reproduce the pipeline end-to-end from public data. Code,
trained weights, per-epoch logs, all metric JSONs, and the six figures are
released under MIT, alongside a public interactive demo.

\section{Experiments \& Results}
All numbers are on held-out cells with a shuffle or identity control.

\begin{table}[h]\centering\small
\begin{tabular}{@{}llcc@{}}
\toprule
Axis & Metric & Model & Control\\
\midrule
Expression recon & $R^2$ & \textbf{0.574} & $-0.574$\\
Spatial (nbr$\to$ctr) & $R^2$ & \textbf{0.331} & $-0.343$\\
Morphology (S$\to$DINOv2) & $R^2$ & \textbf{0.016} & $-0.280$\\
Time transfer & Spearman $\rho$ & \textbf{0.50} & $\sim0$\\
S$\to$image (CFG) & gap @ $w{=}3$ & \textbf{+0.17} & $0$\\
\bottomrule
\end{tabular}
\caption{Every axis beats its control. Morphology $R^2$ is low in absolute
terms because expression$\to$shape is a genuinely weak, many-to-one signal, yet
the gain over shuffle ($+0.30$) is unambiguous.}
\end{table}

\noindent\textbf{Guidance sweep.} Table~\ref{tab:sweep} sweeps the CFG weight
$w$ on 128 held-out cells (epoch-80 checkpoint). Cell-specificity is positive at
\emph{every} $w$ and peaks at $w{=}3$ (gap $+0.173$); it declines by $w{=}5$ as
strong guidance trades condition-fidelity for contrast. Sharpness stays well
above the real-image reference ($0.032$) throughout.

\begin{table}[h]\centering\small
\caption{CFG guidance sweep. gap $=$ MSE-to-shuffled $-$ MSE-to-true; positive
means the image matches its own cell.}
\label{tab:sweep}
\begin{tabular}{@{}cccccc@{}}
\toprule
$w$ & MSE-true & MSE-shuf & gap & sharp & CS\\
\midrule
0 & 0.387 & 0.423 & $+0.036$ & 0.073 & \checkmark\\
1 & 0.368 & 0.515 & $+0.147$ & 0.074 & \checkmark\\
2 & 0.351 & 0.514 & $+0.164$ & 0.070 & \checkmark\\
\textbf{3} & \textbf{0.334} & \textbf{0.507} & $\mathbf{+0.173}$ & \textbf{0.074} & \textbf{\checkmark}\\
5 & 0.401 & 0.522 & $+0.121$ & 0.085 & \checkmark\\
\bottomrule
\end{tabular}\\[2pt]
{\footnotesize CS = cell-specific (gap $>0$).}
\end{table}

\noindent\textbf{Generation.} Figure~\ref{fig:gallery} shows real Xenium cells
(top) beside images generated purely from $S$ (bottom). Bright nuclei, boundaries,
and texture track the real cell; cells with weak expression signal remain
softer---consistent with the weak-signal caveat above.

\noindent\textbf{The decoder evolution, quantified.} Table~\ref{tab:evo} makes
the design history (\S\ref{sec:history}) concrete: each decoder generation fixed
a measurable failure of the last. The MSE decoder was cell-specific but blurry;
the plain diffusion decoder was sharp but hallucinated (gap $\approx0$); CFG is
both sharp and cell-specific.

\begin{table}[h]\centering\small
\caption{Why the decoder changed. ``gap'' is the cell-specificity gap;
``sharp'' the mean gradient magnitude (real $=0.032$).}
\label{tab:evo}
\begin{tabular}{@{}lccl@{}}
\toprule
Decoder & sharp & gap & verdict\\
\midrule
MSE regression & 0.010 & $+$ & correct but blurry\\
Diffusion (no CFG) & 0.053 & $\approx0$ & sharp, hallucinated\\
\textbf{CFG diffusion} & \textbf{0.074} & \textbf{+0.17} & sharp \& cell-specific\\
\bottomrule
\end{tabular}
\end{table}

\noindent\textbf{Time.} In $S$-space the EMT cells trace a monotone trajectory:
distance from the day-0 centroid grows $0\!\to\!2.5\!\to\!4.8\!\to\!5.6\!\to\!6.6$
across $0$d--$7$d, with $\rho{=}0.50$ ($p{<}10^{-190}$) versus a shuffle
$\rho\approx0$. A snapshot-trained state space recovered temporal order in data
it never saw (Fig.~\ref{fig:time}).

\begin{figure}[t]\centering
\includegraphics[width=\linewidth]{diffS_gallery.png}
\caption{Real (top) vs.\ generated-from-$S$ (bottom) cell morphology,
epoch-80 checkpoint.}\label{fig:gallery}
\end{figure}

\begin{figure}[t]\centering
\includegraphics[width=\linewidth]{unified_time_axis.png}
\caption{Cross-platform time transfer: EMT cells move monotonically through the
frozen $S$-space.}\label{fig:time}
\end{figure}

\begin{figure*}[t]\centering
\includegraphics[width=0.96\textwidth]{frontier_comparison.png}
\caption{Frontier-level evaluation. (A) All four axes beat their shuffle
control. (B) Generated-image sharpness ($0.074$) exceeds real cells ($0.032$)---
the output is not a blurry blob. (C) Component-level positioning against
spatial-generative work: the individual pieces exist elsewhere, but the
four-axis fusion with a negative control at every step is ours.}\label{fig:frontier}
\end{figure*}

\section{Impact \& Significance}
The immediate result is modest in scale; the \emph{idea} is not. If a single
learned state can be decoded to any modality, then a dataset that measured only
one axis can have its missing axes \emph{reconstructed}---a spatial-only cohort
could be given plausible morphology, an expression-only atlas a spatial context.
Our time-transfer experiment is a direct, falsifiable demonstration of this
bridge: structure learned in one platform (static Xenium) transferred, with no
retraining, to another (a dynamic EMT course). That is the mechanism by which
fragmented single-cell data---today siloed by assay---could be stitched into one
navigable space. We regard the honest, control-at-every-step methodology as part
of the contribution: the field is accumulating generative models whose sharpness
hides whether they respect their condition, and the cell-specificity gap is a
cheap, general antidote.

\section{Honest Limitations}
(1) Scale: 200k cells, one tissue---below frontier scale. (2) Morphology is a
genuinely weak signal; some generated cells remain blob-like. (3) Correlation,
not causation. (4) Time is a transfer, not native 4D training data. (5) We report
a cell-specificity gap rather than an absolute FID. (6) The final diffusion run
was stopped at epoch 80/120 after the loss had converged ($\sim0.0233$); all
figures use that checkpoint.

\section{Conclusion}
We turned a schematic into a running, servable system: one learned state $S$ that
fuses expression, space, morphology, and (by transfer) time, that generates
cell-specific morphology, and that is validated against an honest control at
every step. The frontier gap is scale and native 4D data---not the idea.

\vspace{0.4em}
{\footnotesize\textbf{Open source \& demo.} All code, trained weights, figures,
and an interactive public demo were produced during the event on public data
(10x Xenium CRC; GSE147405), released under the MIT license. The interactive demo---pick a held-out cell,
generate its morphology from $S$ at an adjustable guidance weight---is public at
\texttt{alexlee-{}-cell-world-model-demo-web.modal.run} (Modal, scale-to-zero
A10G).}

\end{document}
